
\SGARefTarget{sga3:viii:section:7}\section*{7. Appendix: On Monomorphisms of
Group Preschemes}
\addcontentsline{toc}{section}{7. Appendix: On Monomorphisms of Group Preschemes}

The result proved in \SGARefLink{sga3:viii:section:7}{this section} is not needed in the remainder of the
seminar, except in X, \SGARefLink{sga3:x:theorem:8-8}{8.8} and Expos\'es XV and XVI. It depends essentially
on the existence theorem for quotient groups in Expos\'e VI\(_A\).

\SGARefLink{sga3:viii:corollary:5:7}{Corollary 5.7} raises the question of when a monomorphism
\(u:G\to H\) of \(S\)-groups can be asserted to be an immersion, or even a
closed immersion.

We saw in VI\(_B\), \SGARefLink{cor:sga3-VIB-1.4.2}{1.4.2} that this is so when \(S\) is the spectrum of a
field \(k\), provided that \(G\) is of finite type over \(k\) and \(H\) is
locally of finite type over \(k\). It follows readily that the same result
remains valid when \(S\) is merely assumed Artinian.

On the other hand, examples of bijective monomorphisms that are not
immersions are easy to give, for instance over the affine line over a field
or over the spectrum of a discrete valuation ring. Take
\[
  H=(\mathbf Z/2\mathbf Z)_S,\qquad G=G_1\times_SG_2,
\]
where \(G_1\) is the open subgroup of \(H\) complementary to the closed
point \(x\ne0\) in the fiber \(H_s\simeq\mathbf Z/2\mathbf Z\), with \(s\)
a fixed closed point of \(S\), and \(G_2\) is the closed subscheme of \(H\)
which is the union of the subscheme given by the identity section and the
reduced closed subscheme defined by \(x\). One checks immediately that
\(G_2\) is stable under multiplication in \(H\), and hence is a group
scheme. The immersions \(G_i\to H\), \(i=1,2\), then define a homomorphism
of \(S\)-groups
\[
  G=G_1\times_SG_2\longrightarrow H
\]
which is visibly a bijective monomorphism, and even a local immersion, but
is not an immersion. Indeed, \(G\) and \(H\) are reduced; \(G\) has three
disjoint irreducible components, whereas \(H\) has only two.

Notice also that \(G_1\to H\) is an example of an open subgroup of \(H\)
that is not closed, in contrast to the situation for algebraic groups over
a field. The theory of degenerations of elliptic curves gives further
examples of this phenomenon in which \(H\) is smooth over \(S\) of relative
dimension one and \(G\) has connected fibers.

One might hope\begingroup
\renewcommand{\thefootnote}{(\(\ast\))}%
\footnote[1]{In fact, M. Raynaud constructed a counterexample with \(G\)
smooth and with connected fibers; cf. XVI, \SGARefLink{sga3:xvi:subsection:1:1:item:c}{1.1(c)}. Without the connected
fiber hypothesis, one may take
\(S=\operatorname{Spec}(\mathbf Z_2)\), let
\(G=(\mathbf Z/2\mathbf Z)_S\) with its nonidentity closed point removed,
and let \(H=(\mu_2)_S\).}%
\endgroup
that once \(G\) is assumed flat over \(S\), and, say, \(G\) and \(H\) are
of finite presentation over \(S\), every monomorphism \(u:G\to H\) of
\(S\)-groups is automatically an immersion. We shall prove a result of this
kind under additional hypotheses.

First, we may assume that \(S\) is affine and, by the finite-presentation
hypothesis on \(G\) and \(H\), that \(S\) is Noetherian: one reduces to the
spectrum of a ring of finite type over \(\mathbf Z\). Then \(G\) and
\(H\) are Noetherian. Saying that \(u:G\to H\) is a closed immersion
(respectively, an immersion) is then equivalent to saying that \(u\) is a
monomorphism---which holds by hypothesis---and is proper (respectively,
proper at every point of \(u(G)\)).

The valuative criterion for properness shows that it is enough to check,
after every base change \(S'\to S\) with \(S'\) the spectrum of a discrete
valuation ring, complete if desired, that the morphism
\(u':G'\to H'\) has the corresponding properness property. The case of
local properness was omitted from EGA~II, 7.3 and was to appear as an
erratum in EGA~IV.\begingroup
\renewcommand{\thefootnote}{(\(\ast\))}%
\footnote[0]{Cf. EGA~IV\(_3\), 15.7.}%
\endgroup
We are therefore reduced to the case in which \(S\) itself is the spectrum
of a complete discrete valuation ring, subject to the proviso that any
additional hypotheses imposed on \(G,H,u\) be stable under base change.

Let \(s\), respectively \(s'\), be the closed, respectively generic, point
of \(S\). The induced homomorphisms on the fibers,
\[
  u_s:G_s\longrightarrow H_s,
  \qquad
  u_{s'}:G_{s'}\longrightarrow H_{s'},
\]
are closed immersions because they are monomorphisms of group schemes of
finite type over fields (VI\(_B\), \SGARefLink{cor:sga3-VIB-1.4.2}{1.4.2}). We may therefore identify
\(G_{s'}\) with a closed subscheme of \(H_{s'}\). We shall use \SGARefLink{sga3:viii:lemma:7:1}{the
following result}.

\medskip
\noindent\SGARefTarget{sga3:viii:lemma:7:1}\textbf{Lemma 7.1.}
\textit{Let \(S\) be the spectrum of a discrete valuation ring, with generic
point \(s'\). Let \(H\) be an \(S\)-prescheme and let \(L'\) be a closed
subprescheme of the generic fiber \(H_{s'}\), hence also a subprescheme of
\(H\). Its schematic closure \(L\) in \(H\)---the smallest closed
subprescheme of \(H\) containing \(L'\), cf. EGA~I, 9.5---exists. It is
also the unique closed subprescheme of \(H\), flat over \(S\), whose generic
fiber is \(L'\). Formation of \(L\) is functorial in the pair \((H,L')\)
and commutes with fiber products over \(S\). In particular, if \(H\) is an
\(S\)-group and \(L'\) is a subgroup of \(H_{s'}\), then \(L\) is a
subgroup of \(H\).}

The proof is immediate and is left to the reader.\begingroup
\renewcommand{\thefootnote}{(\(\ast\ast\))}%
\footnote[0]{Cf. EGA~IV\(_2\), 2.8.}%
\endgroup
Apply this to the pair \((H,G_{s'})\). The monomorphism \(u:G\to H\)
factors as
\[
  G\xrightarrow{\,u'\,}L\longrightarrow H,
\]
where \(L\) is a subgroup of \(H\), a closed subprescheme flat over \(S\),
and \(u':G\to L\) induces an isomorphism on generic fibers. Thus \(u:G\to H\)
is an immersion (respectively, a closed immersion) if and only if
\(u':G\to L\) is one. We are reduced to the case in which \(H\) is flat over
\(S\) and \(u_{s'}:G_{s'}\to H_{s'}\) is an isomorphism, again subject
to the requirement that any additional hypotheses remain valid when \(H\)
is replaced by a closed subgroup prescheme.

Since \(H\) is then the schematic closure of \(H_{s'}\), if \(u\) is an
immersion its image \(u(G)\) is a subprescheme of \(H\) containing \(H_{s'}\). Its
schematic closure is therefore \(H\), so \(u(G)\) is an open subprescheme of
\(H\). We must consequently prove that \(u\) is an open immersion, or an
isomorphism if the desired conclusion is that it is a closed immersion.

Because \(G\) and \(H\) are flat over \(S\), this is equivalent to the
induced maps on the fibers being open immersions, respectively
isomorphisms (cf. SGA~1, I, \SGARefLink{sga3:viii:corollary:5:7}{5.7}). It already holds for \(u_{s'}\), so it
remains to prove that \(u_s\) is an open immersion, respectively an
isomorphism.

Since \(G\) and \(H\) are flat over \(S\), their fiber dimensions are
constant (VI\(_B\), \SGARefLink{sga3:viii:corollary:4:3}{4.3}). Their generic fibers agree, hence these dimensions
are the same for \(G\) and \(H\). Thus \(u_s:G_s\to H_s\) is a
monomorphism of algebraic groups of the same dimension. It follows that
\(G_s\) is open in \(H_s\), and set-theoretically is a union of connected
components of \(H_s\). Indeed, after passing to an algebraically closed
base field and reducing \(G\) and \(H\), both are smooth over the field and
the assertion is immediate.

The complement \(H_s-G_s\) is closed in \(H_s\), hence in \(H\). Its
complement \(H'\) in \(H\) is open and is stable under the group law. Upon
replacing \(H\) by \(H'\), we may therefore reduce, with the usual proviso,
to the case in which \(u\) is bijective, that is, \(u_s:G_s\to H_s\) is
bijective. In every case we are reduced to proving that \(u\), equivalently
\(u_s\), is an isomorphism, possibly under this bijectivity hypothesis.

Assume first that \(u\) is bijective. If \(H_s\) is reduced, then \(u_s\) is
an isomorphism, because \(G_s\) identifies with a closed subscheme of
\(H_s\) having the same underlying set. In particular, if
\(k=\kappa(s)\) has characteristic zero, every algebraic group over \(k\)
is reduced by Cartier's theorem (cf. VI\(_B\), 1.6.1; VII\(_B\), 3.3.1;
or EGA~IV\(_4\), 16.12.2 and 17.12.5). We have proved:

\medskip
\noindent\SGARefTarget{sga3:viii:proposition:7:2}\textbf{Proposition 7.2.}
\textit{Let \(u:G\to H\) be a homomorphism of group preschemes of finite
presentation over \(S\). If \(u\) is a monomorphism, \(G\) is flat over
\(S\), and every residue field of \(S\) has characteristic zero, then \(u\)
is an immersion.}

Suppose now that \(\kappa(s)\) has characteristic \(p>0\), and restrict to
the case in which \(G\) is commutative. Under the preceding reductions,
\(H\) is also commutative, because it is the schematic closure of
\(H_{s'}\simeq G_{s'}\). For \(n\geq0\), set
\[
  S_n=\operatorname{Spec}(V/\mathfrak m^{n+1}),\qquad
  G_n=G\times_SS_n,\qquad H_n=H\times_SS_n,
\]
where \(V\) is the valuation ring defining \(S\) and \(\mathfrak m\) its
maximal ideal. For \(m>0\), denote by \({}_mG\) and \({}_mH\) the kernels
of multiplication by \(m\) in \(G\) and \(H\). Define similarly
\[
  {}_m(G_n)=({}_mG)_n,
\]
which we abbreviate to \({}_mG_n\), and likewise for \(H\).

By VI\(_A\), \SGARefLink{sga3:viii:corollary:3:2}{3.2}, the quotients
\[
  Q_n=H_n/G_n
\]
exist. Each \(Q_n\) is a commutative group scheme over \(S_n\), flat over
\(S_n\), and \(H_n\to Q_n\) is faithfully flat with kernel \(G_n\).
Formation of the quotient commutes with base extension, so
\[
  Q_n\simeq Q_m\times_{S_m}S_n\qquad(m\geq n).
\]
In particular, \(Q_n\times_{S_n}S_0=Q_0=H_0/G_0\). Since \(G_0\) and
\(H_0\) have the same underlying set, \(Q_0\) is set-theoretically a single
point: it is a purely infinitesimal group.

It follows that every \(Q_n\) is finite and flat over \(S_n\). Let it be
defined by an algebra \(C_n\) over \(V_n=V/\mathfrak m^{n+1}\), finite free
as a \(V_n\)-module. For \(m\geq n\),
\[
  C_n=C_m\otimes_{V_m}V_n,
\]
compatibly also with the comultiplication maps. Hence
\[
  C=\varprojlim C_n
  \quad\text{is a finite free module over}\quad
  V=\varprojlim V_n,
\]
and the comultiplications on the \(C_n\) induce one on \(C\). Thus
\(Q=\operatorname{Spec}(C)\) is a finite flat group scheme over \(S\)
satisfying
\[
  Q\times_SS_n=Q_n
\]
for every \(n\).

\medskip
\noindent\SGARefTarget{sga3:viii:lemma:7:3}\textbf{Lemma 7.3.}
\textit{Let \(K\) be a field and let \(Q\) be a finite group scheme over
\(K\), of degree \(n\). Then the \(n\)-th power morphism on \(Q\) is zero.}

See VII\(_A\), \SGARefLink{sga3:VIIA:8.5}{8.5}.

\medskip
\noindent\SGARefTarget{sga3:viii:remark:7:3:1}\textbf{Remark 7.3.1.}
The statement of Lemma~\SGARefLink{sga3:viii:lemma:7:3}{7.3} still makes
sense for a finite locally free group scheme \(Q/S\) over an arbitrary base
prescheme \(S\). It would be interesting to find a proof in that generality.

Lemma~\SGARefLink{sga3:viii:lemma:7:3}{7.3} (that is, VII\(_A\), 8.5)
does show that the proposed statement holds whenever \(S\) is reduced: apply
it to the fibers of \(Q\) at the maximal points of \(S\), equivalently the
generic points of its irreducible components.

In the preceding situation, where \(S\) is the spectrum of a discrete
valuation ring and \(Q\) is commutative, it follows that
\(n\operatorname{id}_Q=0\). Here \(n\) is a power
\(p^{\nu_0}\) of the residue characteristic. Thus:

\medskip
\noindent\SGARefTarget{sga3:viii:corollary:7:4}\textbf{Corollary 7.4.}
\textit{With \(Q\) and the \(Q_n\) as above, suppose that their common
degree is \(p^{\nu_0}\). Then}
\[
  p^{\nu_0}\operatorname{id}_Q=0,
  \qquad
  p^{\nu_0}\operatorname{id}_{Q_n}=0
  \quad\textit{for every \(n\)}.
\]
